\chapter{Metodologia}
\label{chap:metodologia}

\section{Visão geral da metodologia}
\label{sec:workflow}
\todo{Inserir o workflow completo: geometria $\rightarrow$ módulos aerodinâmico/estrutural/aeroacústico $\rightarrow$ objetivos e restrições $\rightarrow$ otimizador $\rightarrow$ nova geometria.}

\section{Definição do problema de projeto}
\label{sec:definicao_problema}

\subsection{Parâmetros de entrada e condições operacionais}
\todo{Consolidar limitações geométricas, material, número de pás, dados do motor, velocidade de avanço, rotação e condições atmosféricas.}

\subsection{Variáveis de projeto}
\todo{Definir formalmente os parâmetros efetivamente manipulados pelo otimizador: por exemplo, distribuições de corda e torção e demais variáveis que forem mantidas no escopo final.}

\subsection{Funções objetivo e restrições}
\todo{Formalizar matematicamente objetivos aerodinâmicos/acústicos e restrições estruturais e operacionais.}

\section{Módulo aerodinâmico}
\label{sec:met_bemt}
\todo{Descrever a implementação da BEMT: discretização radial, cálculo iterativo, polares, convergência e outputs $T$, $Q$, potência e eficiência.}

\section{Módulo estrutural}
\label{sec:met_estrutural}
\todo{Descrever o modelo estrutural analítico utilizado dentro do loop de otimização.}

\section{Módulo aeroacústico}
\label{sec:met_hanson}
\todo{Descrever a implementação da formulação de Hanson e o indicador acústico escolhido para a função objetivo.}

\section{Integração dos módulos}
\label{sec:integracao}
\todo{Definir interfaces entre os módulos, unidades, estruturas de dados e sequência de execução.}

\section{Algoritmo de otimização}
\label{sec:met_otimizacao}
\todo{Apresentar o método escolhido, população/pontos iniciais, critérios de parada, tratamento de restrições e reprodutibilidade.}

\section{Estratégia de validação}
\label{sec:validacao}

\subsection{Validação aerodinâmica}
\todo{Comparação com dados experimentais de empuxo, torque e/ou potência.}

\subsection{Validação aeroacústica}
\todo{Comparação com medições de SPL e harmônicos de BPF.}

\subsection{Validação estrutural por FEM}
\todo{Usar o FEM como referência de validação do modelo estrutural simplificado, não necessariamente dentro do loop do otimizador. Incluir análise modal se mantida no escopo final.}
